\documentclass[spanish,letterpaper,oneside]{article}

\def\documenttitle {Instrumento de Autoevaluación del Trabajo Final (Módulo 50175)}
\def\documentsubtitle {Fundamentos para la enseñanza y el aprendizaje I · Subproducto 8A (SD26)}
\def\documentsubject {Subproducto Independiente 8A · Autoevaluación de 20 Ítems y Reflexión}
\def\documentauthor {Mtro. Martín Jonathan de la Cruz Muñoz}
\def\coursename {Fundamentos para la enseñanza y el aprendizaje I}
\def\coursecode {MAEM-SD26}

\def\universityname {Instituto de Investigación, Innovación y Estudios de Posgrado para la Educación del Estado de Nuevo León}
\def\universityfaculty {Maestría en Aprendizaje y Enseñanza de las Matemáticas}
\def\universitydepartment {Fundamentos para la enseñanza y el aprendizaje I}
\def\universitydepartmentimage {img/iiiepe/marca-iiiepe.png}
\def\universitydepartmentimagecfg {height=1.75cm}
\def\universitylocation {Monterrey, Nuevo León}
\def\studentfullname {Mtro. Martín Jonathan de la Cruz Muñoz}
\def\studentid {00874}
\def\studentemail {martin.delacruz@campus.iiiepe.edu.mx}
\def\programperiod {2do periodo, septiembre--diciembre de 2026 (SD26)}
\def\programmodality {Modalidad presencial}
\def\programshort {MAEM}

\def\authortable {
  \begin{tabular}{ll}
    \textbf{Alumno:} & \studentfullname \\
    Matrícula: & \studentid \\
    Correo institucional: & \studentemail \\
    Unidad: & \coursename \\
    Programa: & \programshort \\
    Periodo: & \programperiod \\
    & \\
    \multicolumn{2}{l}{Fecha de realización: 19 de noviembre de 2026} \\
    \multicolumn{2}{l}{\universitylocation}
  \end{tabular}
}

\input{template}

\usepackage{helvet}
\renewcommand{\familydefault}{\sfdefault}
\usepackage{array}
\usepackage{longtable}
\usepackage{pdflscape}
\usepackage{tikz}
\usepackage{pgfplots}
\pgfplotsset{compat=1.18}
\usepackage[most]{tcolorbox}
\usetikzlibrary{arrows.meta,positioning,calc,fit,backgrounds,shapes.geometric}

\definecolor{iiiepeNavy}{HTML}{163B5C}
\definecolor{iiiepeBlue}{HTML}{1F6F9F}
\definecolor{iiiepeAqua}{HTML}{20B7A8}
\definecolor{iiiepeOrange}{HTML}{F28C28}
\definecolor{iiiepePaper}{HTML}{F6F8F7}
\definecolor{iiiepeInk}{HTML}{1E2A32}
\definecolor{iiiepePale}{HTML}{EEF5F8}
\definecolor{iiiepeMint}{HTML}{E7F6F4}
\definecolor{iiiepeWarm}{HTML}{FFF3E6}
\def\portraittitlecolor {iiiepeNavy}
\setcitestyle{authoryear,open={(},close={)}}
\makeatletter
\renewcommand{\@biblabel}[1]{}
\makeatother
\setlength{\bibhang}{1.27cm}
\setlength{\bibsep}{6pt}
\AtBeginDocument{\renewcommand{\refname}{Referencias}}

\newcommand{\insertcoverwatermark}{%
  \AddToShipoutPictureBG*{%
    \begin{tikzpicture}[remember picture,overlay]
      \node[opacity=0.04,inner sep=0pt] at (current page.center) {%
        \includegraphics[width=0.65\paperwidth]{img/iiiepe/marca-iiiepe.png}%
      };
    \end{tikzpicture}%
  }%
}

\setlength{\parskip}{5pt}\setlength{\parindent}{0pt}
\begin{document}\insertcoverwatermark\templatePortrait\templatePagecfg\templateFinalcfg\fancyfoot{}\fancyfoot[L]{\scriptsize IIIEPE · MAEM · SD26}\fancyfoot[C]{\scriptsize\thepage}\fancyfoot[R]{\scriptsize Subproducto 8A · Autoevaluación 50175}\setstretch{1.3}\raggedbottom
\section{Propósito del instrumento oficial de autoevaluación (Módulo 50175)}
En congruencia con la guía sumativa del proyecto integrador \textit{Calibrar para decidir en la formación técnica} (Actividad 8, módulo \texttt{50172}), este documento independiente reúne y verifica el \textbf{Instrumento de autoevaluación oficial del módulo \texttt{50175}} para la Maestría en Aprendizaje y Enseñanza de las Matemáticas.

Mi postura es que la autoevaluación docente solo adquiere legitimidad académica cuando cada afirmación marcada responde a evidencia verificable dentro del trabajo entregado y cuando la reflexión personal explicita con honestidad tanto los logros consolidados como los ajustes indispensables antes de implementar el proyecto en grupos reales \citep{bostrom2023formative,restrepo2004investigacion}.

\section{Lista de verificación autoevaluada (20 afirmaciones del módulo 50175)}
A continuación marco el nivel de cumplimiento alcanzado en cada uno de los veinte reactivos institucionales —distinguiendo con precisión la evidencia localizada tanto en la presentación de 15 diapositivas (`.pptx` y Beamer) como en el informe integrador final:

{\small\setstretch{1.12}\renewcommand{\arraystretch}{1.24}
\begin{longtable}{|>{\centering\arraybackslash}p{0.05\linewidth}|>{\raggedright\arraybackslash}p{0.52\linewidth}|>{\centering\arraybackslash}p{0.06\linewidth}|>{\centering\arraybackslash}p{0.08\linewidth}|>{\centering\arraybackslash}p{0.05\linewidth}|>{\raggedright\arraybackslash}p{0.20\linewidth}|}\hline
\rowcolor{iiiepePale} \textbf{\#} & \textbf{Afirmación oficial (Módulo 50175)} & \textbf{Sí} & \textbf{En proc.} & \textbf{No} & \textbf{Evidencia localizada} \\\hline
\endfirsthead
\hline
\rowcolor{iiiepePale} \textbf{\#} & \textbf{Afirmación oficial (Módulo 50175)} & \textbf{Sí} & \textbf{En proc.} & \textbf{No} & \textbf{Evidencia localizada} \\\hline
\endhead
1 & Elaboré una portada con el título del proyecto, mi nombre, la asignatura y la fecha. & X & & & Portada IIIEPE y Diap. 1. \\\hline
2 & Expliqué con claridad en qué consiste mi proyecto. & X & & & Secc. 2 / Diap. 2. \\\hline
3 & Mencioné el nivel educativo al que está dirigido. & X & & & CONALEP/CECyTE e Ing. \\\hline
4 & Describí el propósito del proyecto. & X & & & Meta formativa explícita. \\\hline
5 & Identifiqué los contenidos matemáticos que desarrolla. & X & & & $y=mx+b$ vs. $y=kx$. \\\hline
6 & Describí la problemática o necesidad que atiende mi proyecto. & X & & & Secc. 3 / Diap. 3. \\\hline
7 & Justifiqué por qué esta problemática es relevante. & X & & & Criterios de error y perfil. \\\hline
8 & Expliqué por qué es importante abordarla mediante el aprendizaje de las matemáticas. & X & & & Toma de decisiones reales. \\\hline
9 & Expliqué cómo mi proyecto desarrolla la resolución de problemas. & X & & & Desafío de calibración. \\\hline
10 & Expliqué cómo favorece el pensamiento lógico y matemático. & X & & & Análisis de equivalencia. \\\hline
11 & Expliqué cómo promueve el razonamiento y la argumentación. & X & & & Dictamen técnico por equipo. \\\hline
12 & Expliqué cómo favorece la modelación y representación matemática. & X & & & Tabla, gráfica y ecuación. \\\hline
13 & Relacioné las matemáticas con situaciones de la vida cotidiana. & X & & & Medición y tarifas reales. \\\hline
14 & Identifiqué la metodología utilizada y justifiqué mi elección. & X & & & Indagación STEAM (NEM). \\\hline
15 & Expliqué cómo esta metodología favorece el aprendizaje de las matemáticas. & X & & & 5 fases cíclicas NEM. \\\hline
16 & Describí las actividades principales del proyecto. & X & & & Apertura, núcleo y cierre. \\\hline
17 & Expliqué la relación del proyecto con los sistemas de Bronfenbrenner. & X & & & 5 sistemas ecológicos. \\\hline
18 & Expliqué los aprendizajes matemáticos esperados y su impacto. & X & & & Secc. 8 / Diap. 11--12. \\\hline
19 & Respondí de manera crítica las preguntas de reflexión final. & X & & & 4 respuestas integradas. \\\hline
20 & Incluí al menos tres referencias en formato APA 7. & X & & & Más de 12 referencias APA 7. \\\hline
\end{longtable}}

\section{Reflexión personal del instrumento oficial (Módulo 50175)}
\begin{tcolorbox}[colback=iiiepeMint,colframe=iiiepeAqua!80!black,title=\textbf{Respuestas argumentadas a las dos preguntas de reflexión personal},breakable]
\textbf{1. ¿Qué fue lo que mejor logré en esta actividad?}\\
Logré articular orgánicamente mi trayectoria técnica en mantenimiento industrial, mecatrónica e ingenierías con los marcos pedagógicos estudiados a lo largo de la cohorte SD26 \citep{sawyer2008optimizing,quiroga2022aprendizaje,bronfenbrenner2006bioecologico,sep2022sugerencias} dentro de una sola secuencia realista de modelación matemática de calibración ($y = 1.05x + 0.2$). Además, cuidé explícitamente la honestidad epistémica al distinguir este diseño \textit{por realizar} y sus datos didácticos simulados respecto de la evidencia empírica futura que corresponderá documentar al aplicarlo.\\[6pt]
\textbf{2. ¿Qué aspecto necesito mejorar antes de entregar mi trabajo?}\\
Requiero afinar, al momento de su puesta en marcha con grupos reales, la estimación temporal destinada al trabajo en equipos heterogéneos y sistematizar desde la primera sesión las respuestas diagnósticas iniciales para graduar con mayor precisión las tarjetas de andamiaje escrito según las condiciones laborales y de conectividad de cada plantel \citep{restrepo2004investigacion,eef2021feedback}.
\end{tcolorbox}

\clearpage
\section{Conclusiones}
En conclusión, sostengo que someter el proyecto integrador al escrutinio de estos veinte indicadores fortalece tanto la coherencia didáctica como la responsabilidad ética docente. Por tanto, considero que la articulación entre evidencia conceptual, metodología sociocrítica de la NEM \citep{sep2022plan}, mirada bioecológica \citep{bronfenbrenner2006bioecologico} y autoevaluación crítica deja la propuesta completamente preparada para su revisión institucional y futura implementación en aula.

\clearpage
\begin{thebibliography}{99}
\bibitem[Boström y Palm(2023)]{bostrom2023formative} Boström, E., y Palm, T. (2023). Formative assessment practice in mathematics. \textit{Frontiers in Education}, 8, Artículo 1101192.
\bibitem[Bronfenbrenner y Morris(2006)]{bronfenbrenner2006bioecologico} Bronfenbrenner, U., y Morris, P. A. (2006). The bioecological model of human development. En \textit{Handbook of child psychology} (pp. 793--828). Wiley.
\bibitem[Collin y Quigley(2021)]{eef2021feedback} Collin, J., y Quigley, A. (2021). \textit{Teacher feedback to improve pupil learning}. EEF.
\bibitem[Quiroga Aguilar y Lara Órdenes(2022)]{quiroga2022aprendizaje} Quiroga Aguilar, L., y Lara Órdenes, E. (2022). El aprendizaje profundo como herramienta para cambio en la visión de aprendizaje. \textit{Revista Educación las Américas}, 12(1).
\bibitem[Restrepo Gómez(2004)]{restrepo2004investigacion} Restrepo Gómez, B. (2004). La investigación-acción educativa y la construcción de saber pedagógico. \textit{Educación y Educadores}, (7), 45--55.
\bibitem[Sawyer(2008)]{sawyer2008optimizing} Sawyer, R. K. (2008). Optimizing learning: Implications of learning sciences research. OECD Publishing.
\bibitem[SEP(2022a)]{sep2022plan} Secretaría de Educación Pública. (2022). \textit{Plan de Estudio para la educación preescolar, primaria y secundaria}. SEP.
\bibitem[SEP(2022b)]{sep2022sugerencias} Secretaría de Educación Pública. (2022). \textit{Sugerencias metodológicas para el desarrollo de los proyectos educativos}. SEP.
\end{thebibliography}

\end{document}
